\subsection{The integral lattices used in the determinant}
\label{subsec:division-level-pins}

The determinant compares the integral structures at both levels.

Fix the scale, point index and its chosen twist family. Denote
the division field by \(K\), and use the chosen embedding for
the complex comparison. Let \(S_0,S_1\) be the data at the
base and refined levels of this point, respectively. They come from one realization, so both traces act on the same cut in the ambient representation. Each level
has a separate coefficient lattice \(\Lambda_i\), a Gram
matrix \(X_{i,v}\) at every infinite place \(v\), and a
field-linear coordinate isomorphism \(e_i\). The base map
\(e_0\) takes values in the base row and the refined-level map
\(e_1\) takes values in the row for the period indexed by the
sequence. For \(i=0,1\), write \(\Lambda_i^{\mathrm{div}}\)
for the row lattice in the divisibility calculation.

\begin{proposition}[Level and rank compatibility]\label{prop:division-level-compatibility}
For the compatible level realizations just fixed,
\[
    e_i(\Lambda_i)=\Lambda_i^{\mathrm{div}},\qquad
    \langle e_i z,e_i z'\rangle_v=\langle z,z'\rangle_{X_{i,v}}.
\]
The first equality is over \(\mathcal O_K\) and the metric equality
holds at every infinite place. The base coordinate rank is
the dimension \(D_0\) of the base complex period row and equals
the small-level rank at the base tuple on the specified cut.
\end{proposition}

\begin{proof}
\proofstep{The integral and metric identifications.}
The level character acts by pullback
through the common form map of
Section~\ref{par:selected-form-coordinates}. The identities
come from its inverse image of the integral form lattice
and the Gram of those same basis forms. Applying the
inverse coefficient map recovers the integral lattice
at each level. The metric identity extends from basis
pairs by sesquilinearity. The cut projector is conjugated
by this isomorphism, so its range uses the same integral
intersection under either reading.
Indeed, for a lattice \(L\), a projector \(P\) and the coordinate isomorphism \(T\),
\[
  T(L\cap\operatorname{im}P)=T(L)\cap\operatorname{im}(TPT^{-1}).
\]
 These statements use
the arithmetic form comparison on the selected level,
together with the isometric period comparison at the
chosen complex place. The latter alone would allow any
lattice in the complex span.

\proofstep{The base rank.}
The map \(e_0\) identifies the
base coordinate space with the rational base row, and
its chosen basis becomes a complex basis of the base
period row by the frame comparison. The dimension of
the small-level space at the base tuple on the same cut
therefore equals \(D_0\). The small-level equality across
lifts is the theta adjoint and multiplicity calculation
of Section~\ref{subsec:selected-period}. Its rank is the
dimension before replacing the integral congruence module
by the cut quotient. In the determinant below, both
the base-level geometric covolume and the bound that
divides by \(D_0\) read this same rank. The count of the
analytic retained row and the ramified column comparison
use their own projectors and dimensions.
\label{cor:division-level-pins-1}

\end{proof}

\paragraph{Which tuples occur.}

Which tuples occur?

Let \(W_{\mathrm{tw}}\) be the finite set of twisted places
of the base level and \(\mu(w)\) the tuple assigned there.
The finite character construction covers its tuples by
the base tuple and the images of actual places. The
level realization identifies this place map with the
map \(\mu\) used in the base row, under the fixed
bijection of place indices with the places of \(K\).
For an occurring tuple equal to the base tuple use that tuple. Otherwise take a used twisted
place giving the tuple and send this place through the inverse bijection. Under the fixed bijection of places, the two maps give the same tuple.

In particular, with \(t_0\) the base tuple and \(W_{\rm tw}\) the set of used twisted places, every tuple tested in division belongs to
\[
  \{t_0\}\cup\{\mu(w):w\in W_{\rm tw}\}.
\]
 Therefore the rational span of the lift,
its theta-map identity, and the point-period identity,
which are proved on the used tuples, apply to every
tuple of the division package. In particular, for
a tuple assigned to two places each place keeps its
own valuation and its numerator. Only its rational
period coordinates are shared.

\begin{lemma}[The period on a used tuple]\label{lem:common-division-period}
For each sequence and used tuple of the division package, write \(P_{\rm ar}\) for the bilinear period in the restricted column and row coordinates and \(P_{\rm an}\) for its analytic realization at the chosen complex embedding \(\sigma\). On corresponding vectors \(v,w\),
\[
 \sigma(P_{\rm ar}(v,w))=P_{\rm an}(v_\sigma,w_\sigma).
\]
 The comparison between the base and
sharp levels uses this common period with each level's lattice
and rank.
\end{lemma}

\begin{proof}
Use the restricted column identification and the row map of
Section~\ref{par:selected-form-coordinates}. In bases \((c_i)\)
and \((r_a)\) over the coefficient field, the two period
matrices have the same \((i,a)\)-entry by the equality of the
basis readings. Write \(B\) and \(B'\) for the two period maps in these coordinates. Their agreement extends to arbitrary column and row vectors by
\[
  B\!\left(\sum_i x_i c_i,\sum_a y_a r_a\right)=\sum_{i,a}x_i y_a B(c_i,r_a)=\sum_{i,a}x_i y_a B'(c_i,r_a).
\]
 The coefficient field for this
step is the field of the actual reference kernels.

The word shift and its finite valuations use the
isomorphism on the unshifted column lattice before
multiplying by its fractional ideal. Thus one period map for the sequence,
with that shifted column and the row lattice with the prescribed local conditions,
occurs both in the period comparison and in the
Hermitian slope estimate. The level passage compares the
base and sharp slope using the shared tuple map and the
same period and block scalars at every infinite place.
Each slope uses the lattice and rank at its own level.
\label{cor:division-level-pins-2}

\end{proof}

\paragraph{Transport of the weighted length.}
For a fixed lift \(j\), write \(\eta_j(w)\) for the
normalized length appearing in the uncut record and
\(\eta^{\rm int}_{j,\mu(w)}(w)\) for the integer length in
the base-lattice calculation. With the above place and
tuple comparison one has, at \(w\in W_{\mathrm{tw}}\),
\[
     \eta^{\rm int}_{j,\mu(w)}(w)=e_w\eta_j(w).
\]
Indeed the two coefficient modules agree through the
form and twist maps, and the definition of the normalized
length divides by the absolute ramification index of
that place. Both the finite level condition and the
coefficient comparison are used here. In particular
twisted places are in the support and have \(\mu(w)\ne\mu_0\).
Their nonnegative contribution is bounded by \(H_j\)
of Section~\ref{subsec:division-length}.

The local numerator calculation uses the marked
representation at each surviving character and the
multiplicity of that character on this cut. The comparison
back to the original module uses the original level
operator and its order. The positive local length table
is uniform over the retained character components at a fixed
point. Its use on the original row needs the separate
rank and character conditions of Sections
\ref{subsec:selected-cut} and~\ref{subsec:division-length}.
On these inputs the bound there gives
\[
  \frac{1}{D_0}
    \sum_{w\in W_{\mathrm{tw}}}\omega_w
                          \eta^{\rm int}_{j,\mu(w)}(w)
        \le (|C_N|+1)\,\mathcal E ,
       \qquad \mathcal E=z(B_0+N_c),\quad z=x^\epsilon .
\]
Here \(C_N\) is the constant of the uncut length bound.
The strict cohomology contributes at most \(\sqrt b\)
per unit of rank. Extension to a local splitting field
multiplies unnormalized length and absolute ramification
by the same factor. The sum over coefficient places
therefore charges each rational prime at most \(\log\ell\)
times its normalized bound. This keeps the strict budget
on the same host and within the same rank cutoff.

The determinant comparison has a further finite cost.
Let \(d_{j,t}\) be the class denominator of the block at
lift \(j\) and tuple \(t\), including the level index in the
class form. Choose its common multiple \(d_*\) over
all lifts and tuples, and let \(s\) be the value-lattice
multiplier. The preceding bound for the uncut numerator gives the
following estimate including the dual-inclusion correction.

\begin{auxlemma}[A common multiple of all stabilizer orders]
\label{lem:uniform-stabilizer-denominator}
For every \(d\ge1\) there is an integer \(M(d)\) divisible by the
order of every finite subgroup of \(\mathrm{GL}_d(\mathbb Z)\), with the absolute bound
\[
                    \log M(d)\ll d\log(2d).
\] In particular
the l.c.m. \(D_\Gamma\) of the finite class stabilizer orders in the
definite row satisfies \(\log D_\Gamma\ll A\log A\) for \(A=fN_m\ge2\).
\end{auxlemma}
\begin{proof}
A finite subgroup injects under reduction modulo an odd prime \(q\).
Indeed, if a nonidentity element in the kernel has prime order \(p\),
write it as \(1+q^h U\), where \(h\ge1\) and \(U\not\equiv0\pmod q\).
The first nonconstant term of its \(p\)-th power has exact order
\(h\) if \(p\ne q\), and \(h+1\) if \(p=q\). All higher terms
have greater order, a contradiction.

Fix an odd prime \(p\). By Dirichlet's theorem choose an odd prime
\(q\) in a primitive root class modulo \(p^2\). Counting ordered bases gives
\[
       |\mathrm{GL}_d(\mathbb F_q)|=q^{d(d-1)/2}\prod_{i=1}^d(q^i-1).
\]
Only the factors \(q^i-1\) contribute to the \(p\)-part. Their orders vanish
unless \(p-1\mid i\), and at \(i=(p-1)j\) are \(1+v_p(j)\) by the
binomial theorem. With \(b_p=\lfloor d/(p-1)\rfloor\), the
\(p\)-valuation of every subgroup order is therefore at most
\[
    b_p+v_p(b_p!)\le 2b_p .
\]

For \(p=2\) take \(q=3\). The estimate
\(v_2(3^i-1)\le2+v_2(i)\) gives a total at most \(3d\).
Thus one may use
\[
 M(d)=2^{3d}\prod_{\substack{3\le p\le d+1\\p\ {\rm prime}}}p^{2b_p},
 \qquad \log M(d)\ll d+d\sum_{p\le d+1}\frac{\log p}{p}
                   \ll d\log(2d).
\]
For the last bound each dyadic interval \((m,2m]\) contributes
\(O(1)\) to the prime sum, since its primes divide
\(\binom{2m}{m}\le4^m\).

Each class stabilizer acts faithfully on its positive definite
Hermitian lattice. Forgetting its \(\mathcal O_E\)-action gives
an integral lattice of rank at most \([E:\mathbb Q]N_m=2A\).
Adjoining a trivial summand embeds each such stabilizer in
\(\mathrm{GL}_{2A}(\mathbb Z)\), so their l.c.m. divides \(M(2A)\).
\end{proof}

\begin{proposition}[Weighted dual-lattice cost]\label{prop:weighted-dual-cost}
For the common level realization of
Proposition~\ref{prop:division-level-compatibility}, each lift \(j\),
and a positive
integer value-lattice multiplier \(s\),
\begin{align*}
 &\frac{1}{D_0}
     \sum_{w\in W_{\mathrm{tw}}}\omega_w
        \bigl(\eta^{\rm int}_{j,\mu(w)}(w)
                      +D_0 v_w(d_{j,\mu(w)}s)\bigr) \\
 &\hspace{40pt}\le (|C_N|+1)\mathcal E+\log d_*+\log s .
\end{align*}
At \(s=1\), the least common multiple \(d_*\) satisfies
\(\log d_*\le(2+9C_\Gamma)\mathcal E\) on the reference size range \(fN_m\le3\sqrt b\), \(\log(fN_m)\le3z\) of \eqref{eq:weighted-dual-reference-classes}.
\end{proposition}

\begin{proof}
\proofstep{Valuations of the common denominator.}
For the correction, divisibility bounds the valuation
at each individual place by \(v_w(d_*s)\).
The sum of the latter over all finite places with
weights \(\omega_w\) is \(\log(d_*s)\) by the rational
norm formula. This argument allows \(t=\mu(w)\) to vary
with \(w\). The sum of the least exponents is handled
in the first error separately.

\proofstep{Size of the denominator in the fixed comparison.}
In the fixed comparison \(s=1\), and \(d_*\) can be the
least common multiple of the block denominators.
Thus for the list of block denominators the integral choice is
\[
  s=1,\qquad d_*=\operatorname{lcm}_b d_b .
\]
Each block denominator divides \(|G|\,D_{\Gamma}\)
for one reference member, where \(G\) is the finite row
group and \(D_{\Gamma}\) is the common multiple of its
class stabilizer orders. Indeed the block level subgroup
order divides \(|G|\) and its class stabilizer orders
divide \(D_{\Gamma}\) by the class correspondence.
Consequently \(d_*\) also divides this product.

In valuations, for block denominators \(d_b\), this is
\[
  v_p(d_*)=\max_b v_p(d_b)\le v_p(|G|)+v_p(D_\Gamma)\qquad\text{for every prime }p.
\]
 The
group-size comparison gives \(\log|G|\le2\mathcal E\).
Each class stabilizer acts faithfully on an integral lattice of rank at most \(2fN_m\). Lemma~\ref{lem:uniform-stabilizer-denominator} gives the first inequality below. The other two are the size bounds for the marked definite row:
\begin{equation}\label{eq:weighted-dual-reference-classes}
 \log D_{\Gamma}\le C_\Gamma A\log A,\qquad
 1\le A=fN_m\le3\sqrt b,\qquad \log A\le3z .
\end{equation} Thus
\(A\log A\le9bz\le9\mathcal E\), and
\(\log d_*\le(2+9C_\Gamma)\mathcal E\). These two
logarithms come from the class dual inclusion and the
value multiplier. The passage from the cut to the
uncut numerator has already been taken before this
comparison.
\label{cor:division-level-pins-3}
\end{proof}
